\documentclass[../../../include/open-logic-section]{subfiles}

\begin{document}

\olfileid{sth}{replacement}{extrinsic}
\olsection[બાહ્ય વિચારણાઓ]{પ્રતિસ્થાપન અંગે બાહ્ય વિચારણાઓ}

શરૂઆતમાં પ્રતિસ્થાપનને વાજબી ઠરાવવાનો \emph{બાહ્ય} પ્રયાસ
વિચારીએ. બૂલોઝ આવો એક પ્રયાસ સૂચવે છે:
\begin{quote}
  [\ldots] પ્રતિસ્થાપનનાં સ્વયંસિદ્ધિઓ સ્વીકારવાનું કારણ એકદમ
  સરળ છે: તેમનાં ઘણાં ઇચ્છનીય પરિણામો છે અને (દેખીતી રીતે)
  કોઈ અનિચ્છનીય પરિણામો નથી. પુનરાવર્તી કલ્પના અંગેના
  પ્રમેયો ઉપરાંત, પરિણામોમાં અનંત સંખ્યાઓનો સંતોષકારક,
  ભલે આદર્શ ન હોય એવો સિદ્ધાંત, તથા સુઆધારિત સંબંધો પર
  અનુમાનાત્મક વ્યાખ્યાઓને વાજબી ઠરાવતું અત્યંત ઇચ્છનીય
  પરિણામ સામેલ છે.
  \citep[229]{Boolos1971}
\end{quote}
બૂલોઝના વિચારનો સાર એ છે કે પ્રતિસ્થાપનને તેનાં ફળોને લીધે
વાજબી ઠરાવવું. તેમણે જણાવેલાં ચોક્કસ ફળો છેલ્લાં થોડાં
પ્રકરણોમાં આપણે ચર્ચેલાં પરિણામો છે. પ્રતિસ્થાપનથી આપણે
સાબિત કરી શક્યા કે વૉન નોયમન ક્રમવાચક સંખ્યાઓ સુક્રમના
ક્રમપ્રકારના ખ્યાલના ઉત્તમ અવેજો છે (આ ``અનંત સંખ્યાઓનો
સંતોષકારક, ભલે આદર્શ ન હોય એવો સિદ્ધાંત'' છે).
પ્રતિસ્થાપનથી આપણે $V_\alpha$ની વ્યાખ્યા આપી, સ્તરાંકનો
ખ્યાલ સ્થાપ્યો અને $\in$-અનુમાન સાબિત કર્યું (આ
``પુનરાવર્તી કલ્પના અંગેના પ્રમેયો'' છે). છેવટે,
પ્રતિસ્થાપનથી અનંતાતીત પુનરાવર્તનનું પ્રમેય સાબિત થાય છે
(આ ``સુઆધારિત સંબંધો પર અનુમાનાત્મક વ્યાખ્યાઓ'' છે).

આ પરિણામો ખરેખર ઇચ્છનીય છે. પરંતુ શું એ પરિણામો
પ્રતિસ્થાપનને \emph{વાજબી ઠરાવવા} પૂરતાં છે? \emph{ના}.
અથવા ઓછામાં ઓછું, સીધી રીતે તો નહીં.

એક સરળ સમસ્યા આ છે. પ્રતિસ્થાપનનાં કેટલાંક ઇચ્છનીય
પરિણામો બતાવ્યાં છે, પરંતુ તેમાંથી ઘણાં આપણે બીજી રીતોથી
પણ મેળવી શક્યા હોત. આ વાત જેટલી જાણીતી હોવી જોઈએ
તેટલી નથી; તેથી થોભીને પરિસ્થિતિ સમજીએ.

ગણોનો એક સરળ સિદ્ધાંત છે, સ્તર સિદ્ધાંત, ટૂંકમાં
$\LT$.\footnote{$\LT$નાં પહેલાનાં સ્વરૂપો \citet{Montague1965}
અને \citet{Scott1974}એ રજૂ કર્યાં; \citet{Potter2004}એ તેને
સરળ બનાવીને પુસ્તક જેટલું વિસ્તૃત નિરૂપણ આપ્યું; અને
તાજેતરમાં \citet{ButtonLT1}એ $\LT$ને વધુ સરળ બનાવ્યો.}
$\LT$નાં સ્વયંસિદ્ધિઓ માત્ર ઘટકો દ્વારા નિર્ધારિત સમાનતા,
પૃથક્કરણ અને
દરેક ગણ કોઈ \emph{સ્તર}નો ઉપગણ છે એવો દાવો છે. અહીં
``સ્તર''ની વ્યાખ્યા કુશળતાથી એવી અપાઈ છે કે સ્તરો
આપણા પરિચિત $V_\alpha$ની જેમ વર્તે. તેથી $\ZF$માંથી
$\LT$ સાબિત થાય છે; પરંતુ $\LT$ એ $\ZF$ કરતાં
\emph{ઘણો} નબળો છે. ખરેખર, $\LT$માંથી જોડ, ઘાતગણો,
અનંતતા કે પ્રતિસ્થાપન મળતાં નથી. $\LT$માં અનંતતા અને
ઘાતગણો ઉમેરીને મળતા સિદ્ધાંતને $\Zr$ કહીએ; તેમાંથી
જોડ પણ મળે છે, એટલે $\Zr$ ઓછામાં ઓછો $\Z$ જેટલો
શક્તિશાળી છે. પરંતુ હકીકતમાં $\Zr$ એ $\Z$ કરતાં
કડક રીતે વધુ શક્તિશાળી છે, કારણ કે તેમાં દરેક ગણને સ્તરાંક
હોય છે એવો દાવો ઉમેરાય છે (એટલે તેને $\Zr$ કહેવાનું
મારું સૂચન). ખરેખર, $\Zr$માંથી ક્રમવાચક સંખ્યાઓનો
પૂરેપૂરો સંતોષકારક સિદ્ધાંત, પદાનુક્રમને સુક્રમિત તબક્કાઓમાં
વહેંચતાં પરિણામો, $\in$-અનુમાનની સાબિતી અને અનંતાતીત
પુનરાવર્તનનું એક \emph{સ્વરૂપ} મળે છે.

ટૂંકમાં, બૂલોઝને આ ખબર ન હોવા છતાં, તેમણે જણાવેલાં
બધાં ઇચ્છનીય પરિણામો પ્રતિસ્થાપન \emph{વિના} મેળવી
શકાયાં હોત; તેમણે $\Z$ને બદલે $\Zr$ વાપરવો પડ્યો હોત.

(આ બધું હોય, તો આપણે $\LT$ અને $\Zr$ને બદલે
$\ZF$ શીખવવાનો રૂઢ માર્ગ શા માટે અપનાવ્યો? બે કારણો છે.
પહેલું, માત્ર ઐતિહાસિક કારણોથી $\LT$થી શરૂઆત કરવી
બહુ પ્રચલિત નથી; અમે તમને ગણસિદ્ધાંતની વધુ પ્રમાણભૂત
ચર્ચાઓ વાંચવા સજ્જ કરવા માગતા હતા. બીજું, તમે $\LT$
અને $\Zr$ સમજવા તૈયાર થશો ત્યારે સીધાં \citealt{Potter2004}
અને \citealt{ButtonLT1} વાંચી શકો છો.)

અલબત્ત, $\Zr$ એ $\ZF$ કરતાં કડક રીતે નબળો હોવાથી,
$\ZF$ સાબિત કરે પણ $\Zr$ ખુલ્લાં રાખે એવાં પરિણામો છે.
આવાં પરિણામોમાંથી કોઈ એક બતાવીને બાહ્ય આધારે
પ્રતિસ્થાપનને વાજબી ઠરાવવાનો પ્રયાસ કરી શકાય. પરંતુ
$\Zr$નો ઉપયોગ શીખ્યા પછી એવી બાબતોનાં વધુ ઉદાહરણો
શોધવાં મુશ્કેલ પડે છે જે (ક) પ્રતિસ્થાપનથી, પણ બીજી
રીતે નહીં, નક્કી થાય અને (ખ) સહજ રીતે સાચી લાગે.
(આ અંગે વધુ માટે \citealt[\S13.2]{Potter2004} જુઓ.)

મુદ્દો આ છે: પ્રતિસ્થાપનનું ખાતરીકારક બાહ્ય ઔચિત્ય
આપવા માટે એવું પરિણામ શોધવું પડે જે પ્રતિસ્થાપન વિના
\emph{મેળવી જ ન શકાય}. એ સહેલું કામ નથી.

હવે પ્રતિસ્થાપન માટેના કોઈ પણ સંપૂર્ણ બાહ્ય ઔચિત્ય-પ્રયાસની
બીજી સમસ્યા જોઈએ. (આ સમસ્યા કદાચ પહેલી કરતાં વધુ પાયાની
છે.) બૂલોઝ એટલું જ કહેતા નથી કે પ્રતિસ્થાપનનાં ઘણાં
ઇચ્છનીય પરિણામો છે. તેઓ એમ પણ કહે છે કે તેનાં
``(દેખીતી રીતે) કોઈ અનિચ્છનીય'' પરિણામો નથી. પરંતુ
કૌંસમાં મુકાયેલી ``દેખીતી રીતે'' એ સાવધાની નિશ્ચિતપણે
અત્યંત મહત્ત્વની છે.

અહીં સુધી આપણે કેવી રીતે પહોંચ્યા તે યાદ કરો: સહજ
ગણરચનાથી અસુસંગતતા ઊભી થઈ, અને તેના જવાબમાં આપણે
ગણની સંચયી-પુનરાવર્તી કલ્પના સ્વીકારી. આ કલ્પના સાથે
એવી એક વાત મળે છે જે, આશા છે કે, તેની સુસંગતતા વિશે
આપણને ખાતરી આપે છે. પરંતુ એ વાતની અંદરથી પ્રતિસ્થાપનને
વાજબી ન ઠરાવી શકીએ, તો $\ZF$ સુસંગત છે એમ માનવાનું
આપણને (હજી સુધી) કોઈ કારણ નથી. વધુ ચોક્કસ રીતે,
હજી સુધી કોઈએ વિરોધાભાસ \emph{શોધ્યો નથી} એ
(કદાચ માત્ર આકસ્મિક) હકીકત સિવાય $\ZF$ સુસંગત છે
એવું માનવાનું કોઈ
કારણ નથી. એ જ અર્થમાં બૂલોઝની ટિપ્પણીનો સાર જાણે
``(દેખીતી રીતે) $\ZF$ સુસંગત છે'' એટલો જ થાય છે.
સુસંગતતા અંગે આપણે આનાથી વધુ ખાતરી માગવી જોઈએ.

પ્રતિસ્થાપનને \emph{સંપૂર્ણપણે} બાહ્ય રીતે વાજબી
ઠરાવવાના, એટલે કે માત્ર $\ZF$નાં (જાણીતાં) પરિણામોના
આધારે અપાતા, કોઈ પણ પ્રયાસને આ સમસ્યા લાગુ પડે છે.
તેથી આપણે પ્રતિસ્થાપનનું \emph{આંતરિક} ઔચિત્ય શોધવા
માગીશું: ગણો વિશે આપણે આપેલી વાત જાણે ``પહેલેથી જ''
પ્રતિસ્થાપન સ્વીકારવા બાંધે છે એવું દર્શાવતું ઔચિત્ય.

\end{document}
